\clearpage
\chapter*{Erratum et complément à l'exposé n\textsuperscript{o} 236}
\addcontentsline{toc}{chapter}{Erratum et complément à l'exposé
n\textsuperscript{o} 236}
\src{303}

\noindent\textbf{8. TDTE VI : Les schémas de Picard. Propriétés générales.}

\begin{center}
[Séminaire Bourbaki, t. 14, 1961/62, n\textsuperscript{o} 236, 23 p.]
\end{center}

\noindent\textbf{Page 12, ligne 16.} --- Au lieu de « droite affine », lire
« droite affine privée de l'origine ».

\noindent\textbf{Ligne 18.} --- Au lieu de « $X[t]$ », lire
« $X[t,t^{-1}]$ ».

\noindent\textbf{Page 16, proposition 3.1, 6\textsuperscript{e} ligne.} ---
Au lieu de « Or l'hypothèse ... », lire « Or, de l'hypothèse ... ».

\noindent\textbf{Page 21, § 4 : Le théorème de finitude pour le schéma de
Picard.} --- Les questions de finitude du genre de celles soulevées dans ce
paragraphe ont été à peu près totalement résolues depuis la rédaction du
présent exposé. Indiquons les faits principaux connus maintenant dans cette
direction. Pour simplifier les énoncés, nous supposons implicitement que tous
les préschémas de Picard intervenant dans les énoncés existent, bien qu'une
modification évidente de ces énoncés permette de se débarrasser de toute
hypothèse d'existence. $S$ désigne par la suite un schéma noethérien, $X,Y$ des
schémas propres sur $S$.

\medskip
\noindent(i) Soit $f:X\longrightarrow Y$ un morphisme surjectif, alors
\[
f^*: \operatorname{Pic}_{Y/S}\longrightarrow\operatorname{Pic}_{X/S}
\]
est un morphisme de type fini.

\src{304}
\noindent(ii) Supposons $Y$ projectif sur $S$, muni d'un Module inversible
ample relativement à $S$, et soit $X$ le schéma des zéros d'une section
quelconque de ce Module, $f:X\longrightarrow Y$ l'immersion canonique.
Supposons enfin que les composantes irréductibles des fibres de $Y/S$ soient de
dimension $\geq3$. Alors
\[
f^*: \operatorname{Pic}_{Y/S}\longrightarrow\operatorname{Pic}_{X/S}
\]
est de type fini.

\noindent(iii) Supposons que $X$ soit projectif sur $S$, et que toutes ses
fibres géométriques soient intègres et de dimension $n$. Soit
$\mathcal O_X(1)$ un Module inversible ample sur $X$, permettant de définir
des polynômes de Hilbert. Soit $M$ une partie de
$\operatorname{Pic}_{X/S}$ ; alors $M$ est quasi-compacte si et seulement si,
dans les polynômes de Hilbert
\[
a_0x^n+a_1x^{n-1}+a_2x^{n-2}+\cdots
\]
des éléments de $M$, les coefficients $a_1$ et $a_2$ restent bornés, en omettant tout coefficient dont le monôme aurait un degré négatif ; ainsi seul $a_1$ intervient si $n=1$, et aucun des deux si $n=0$.

\noindent(iv) Pour tout entier $n\neq0$, l'homomorphisme de puissance
$n$-ième dans le préschéma en groupes $\operatorname{Pic}_{X/S}$ est un
morphisme de type fini.

Notons que (i) et (ii) signifient aussi que (sous les hypothèses envisagées
dans ces énoncés (i) et (ii)) une partie $M$ de
$\operatorname{Pic}_{Y/S}$ est quasi-compacte si et seulement si son image
dans $\operatorname{Pic}_{X/S}$ l'est. On en conclut qu'un Module inversible
$\mathcal L$ sur $Y$ est $\tau$-équivalent à zéro si et seulement si son image
inverse sur $X$ l'est ; en d'autres termes
$\operatorname{Pic}_{Y/S}^\tau$ est l'image inverse de
$\operatorname{Pic}_{X/S}^\tau$. En particulier, pour montrer que le premier
préschéma est de type fini sur $S$, il suffit de le prouver pour le second
puisque $\operatorname{Pic}_{Y/S}\longrightarrow
\operatorname{Pic}_{X/S}$ est de type fini. Utilisant alors (i), le lemme de
Chow et (iii), on trouve :

\medskip
\noindent(v) $\operatorname{Pic}_{X/S}^\tau$ est de type fini sur $S$.

De façon générale, la conjonction de (i) pour un morphisme fini et de (ii)
permet de se ramener, pour la plupart des questions de finitude, au cas où
$X/S$ est à fibres géométriques intègres normales de dimension $\leq2$ ;
souvent même, appliquant (i) pour un morphisme surjectif non nécessairement
fini, et la résolution des singularités des surfaces algébriques (due en
caractéristique quelconque à Abhyankar), on se ramène au cas où $X/S$ est même
simple, donc à fibres géométriques non singulières de dimension 2. Cela permet
par exemple, compte tenu de (v), et de l'inégalité de Picard--Igusa majorant le
rang du groupe de Néron--Severi d'une surface projective non singulière, de
prouver la généralisation suivante du théorème de finitude de Néron :

\medskip
\noindent(vi) Soit $X/S$ propre sur $S$ et par ailleurs quelconque, alors les
groupes de Néron--Severi
\[
\operatorname{Pic}_{X_i/k_i}/\operatorname{Pic}_{X_i/k_i}^\circ
\]
des fibres géométriques $X_i/k_i$ de $X/S$

\src{305}
sont de type fini, et leur rang et l'ordre de leur sous-groupe de torsion
restent majorés.

La même méthode de réduction au cas des surfaces non singulières, et des
théorèmes connus pour ce cas (savoir le théorème de finitude de Néron, et le
fait que la forme intersection sur le groupe de Néron--Severi est non
dégénérée) impliquent :

\medskip
\noindent(vii) Soit $X$ un schéma propre sur un corps algébriquement clos.
Alors il existe un nombre fini de courbes fermées intègres $C_i$
($1\leq i\leq r$) dans $X$, ayant la propriété suivante : pour toute partie
$M$ de $\operatorname{Pic}_{X/k}$, $M$ est quasi-compacte si et seulement si
les entiers
\[
\deg_{C_i'}\mathcal L\qquad(\mathcal L\in M)
\]
restent bornés (où $C_i'$ désigne la normalisée de $C_i$).

On peut prendre ici pour $r$ le rang du groupe de Néron--Severi. Une fois connu
que ce dernier est de type fini, (vii) se réduit au fait que les formes
linéaires sur le groupe de Néron--Severi définies par des courbes $C$ dans $X$
ne s'annulent simultanément que sur les éléments de torsion du groupe de
Néron--Severi. Dans le cas $X$ projectif non singulier, ce résultat, ainsi que
(v), était dû à Matsusaka. Utilisant (vii), on obtient facilement la
caractérisation suivante des Modules inversibles $\tau$-équivalents à 0 sur
$X$ :

\medskip
\noindent(viii) Soient $X/k$ un schéma propre sur un corps, $\mathcal L$ un
Module inversible sur $X$. Conditions équivalentes :

\noindent a. $\mathcal L$ est $\tau$-équivalent à 0 ;

\noindent b. Pour tout Module cohérent $F$, on a
\[
\chi(F\otimes\mathcal L)=\chi(F),
\]
où $\chi$ désigne la caractéristique d'Euler--Poincaré ;

\noindent b'. Comme (b), mais avec $F=\mathcal O_Y$, $Y$ étant un
sous-schéma fermé intègre de dimension 1 dans $X$ ;

\noindent c. Pour tout $Y$ comme ci-dessus, désignant par $Y'$ la courbe
normalisée, on a
\[
\deg\mathcal L_{Y'}=0.
\]

Si $X/k$ est projectif et muni d'un Module inversible ample
$\mathcal O_X(1)$, les conditions précédentes sont aussi équivalentes aux
suivantes :

\noindent d. Pour tout entier $m$, $\mathcal L^{\otimes m}(1)$ est ample.

\noindent e. (Si $X$ est intègre.) Pour tout couple d'entiers $m,n$, on a
\[
\chi(\mathcal L^{\otimes m}(n))=\chi(\mathcal O_X(n))
\]

\src{306}
(i.e. (b) est vrai en faisant $F=\mathcal L^{\otimes m}(n)$).

Pour la suffisance de cette dernière condition, on notera qu'elle signifie que
les polynômes de Hilbert des $\mathcal L^{\otimes m}$ sont tous égaux, donc,
en vertu du critère de Mumford (iii), les $\mathcal L^{\otimes m}$ restent
dans un ensemble quasi-compact de $\operatorname{Pic}_{X/k}$, i.e. on a (a).
Les conditions (b), (b'), (c) doivent être considérées comme des variantes
(sur un schéma propre quelconque) de la notion d'équivalence numérique, définie
habituellement sur des variétés projectives non singulières. Pour ces
dernières, l'équivalence de (a) et (c) était évidemment bien connue
(Matsusaka).

Le critère (e) implique aussi le résultat suivant :

\medskip
\noindent(ix) Soit $f:X\longrightarrow S$ un morphisme projectif et plat dont
les fibres géométriques sont intègres. Alors
$\operatorname{Pic}_{X/S}^\tau$ est ouvert et fermé dans
$\operatorname{Pic}_{X/S}$.

Nous nous bornerons à quelques commentaires sur la démonstration des
résultats-clefs (i), (ii), (iii) (le résultat (iv) est un peu à part des
autres, et se prouve en utilisant seulement (i) pour les morphismes finis
surjectifs radiciels, de façon précise, pour un morphisme de Frobenius). Pour
(i), on utilise de façon essentielle les idées de la descente non plate (voir
notamment TDTE I, page 9). Il se trouve que (n'ayant en vue que des résultats
de finitude), le manque de critères d'effectivité pour des données de descente
est inoffensif. D'autre part, Mumford a démontré récemment une forme un peu
moins forte de (iii), savoir le critère faisant intervenir tous les
coefficients du polynôme de Hilbert. Son argument est extrêmement simple, et
inspiré par la démonstration d'un critère d'amplitude de Nakai (énoncé par ce
dernier pour les surfaces non singulières, et généralisé par Mumford aux
morphismes projectifs quelconques). Il me semble d'ailleurs que cet argument
n'est valable que moyennant une légère restriction supplémentaire sur les
fibres de $X/S$, la propriété $(S_2)$ de Serre, vérifiée si les fibres
géométriques sont normales. On utilise alors ce critère restreint dans la
démonstration de (ii) : le critère (i) permet en effet de se ramener au cas où
$Y/S$ est plat à fibres géométriques intègres et normales, et appliquant le
critère de Mumford, on se ramène aisément au cas où $X/S$ satisfait aux mêmes
conditions. De l'hypothèse de dimension résulte alors que les fibres
géométriques de $Y$ et de $X$ sont de profondeur $\geq2$ en leurs points
fermés, ce qui permet d'appliquer les « critères d'équivalence » sous la forme
qui leur est donnée dans SGA 1962 (exposé XII), et d'achever la démonstration de
(ii). Une fois qu'on dispose de (i) et (ii), il n'est pas difficile dans le
critère de Mumford de se débarrasser de toute hypothèse de normalité sur les
fibres, et de le démontrer également sous la forme plus forte donnée dans
(iii).

\src{307}
Notons enfin que la démonstration de (i) et (ii) montre également que dans le
cas où $S$ est le spectre d'un corps $k$, le morphisme
\[
\operatorname{Pic}_{Y/k}\longrightarrow\operatorname{Pic}_{X/k}
\]
est affine (et non seulement de type fini).

\par\smallskip
\begingroup\footnotesize
\noindent\textit{La lecture imprimée demande que les coefficients $a_1$ et $a_2$ restent bornés, sans préciser le cas $n<2$. La lecture courante explicite seulement la convention des coefficients absents en basse dimension. Il ne s'agit pas d'une assertion imprimée démontrée fausse.}\par
\endgroup
